\chapter{Outline of the proof}\label{s3:chapter}

This section explains the main steps of the proof and the ideas behind them. It is meant to be read before the technical parts, and it fixes the vocabulary used there. We describe the proof of \cref{s2:thm:main} in \cref{s3:strategy,s3:rg,s3:fixed,s3:largefield,s3:flow,s3:stochastic,s3:construction}. The proofs of the results on short distances and on weak coupling are outlined in \cref{s3:short,s3:weak}. Throughout, $L$ and $\gamma$ are fixed, $g_3^2=\gamma/L$, and $M\ge3$ is a fixed odd integer.

\section{Strategy}\label{s3:strategy}

The Wilson measure $\mu_{N,\gamma}$ at fine mesh $a_N$ lives on a space whose dimension grows like $N^3$. The observables we care about, flowed Wilson loops at a fixed positive time $t$, depend on the configuration at length scales of order $\sqrt t$ and larger. The renormalization group organizes the integration over the configuration scale by scale, from the finest lattice to a coarse lattice of fixed side. At each scale, we only need to control a weakly coupled and approximately Gaussian integral. In three dimensions the effective coupling at lattice side $m$ is $g_m^2=\gamma/m$, which is small when $m$ is large. For this reason the renormalization group is a convergent procedure in the ultraviolet, as Bałaban showed \cite{Balaban1985UV}. However, this is true only as long as the field is close to a pure gauge at every scale. Configurations for which this fails, the large fields, have to be treated separately.

We organize the proof as a limit in two stages.
\begin{enumerate}
\item We fix a coarse side $u$ and a large integer $K$, and let $N=uB^K$, where $B=M^J$ is the blocking factor of one renormalization group step. We condition $\mu_{N,\gamma}$ on the small-field event $S_{N,u,\gamma}$ of \cref{s2:largefield}, which gives the conditioned Wilson measure $\nu_{N,u,\gamma}=\mu_{N,\gamma}(\,\cdot\mid S_{N,u,\gamma})$. On this event we can run the renormalization group exactly from side $N$ down to side $u$, without any analysis of large fields. We prove that, for each fixed $u$, the laws of flowed Wilson loops under $\nu_{uB^K,u,\gamma}$ converge as $K\to\infty$.
\item By \cref{s2:thm:large-field}, we have $\mu_{uB^K,\gamma}(S_{uB^K,u,\gamma}^c)\le F(u)$ for all $K$, where $F(u)\to0$ as $u\to\infty$. For a bounded observable $G$ we have the elementary inequality
\begin{equation}\label{s3:eq:transfer}
 \bigl\lvert\E_{\mu_{N,\gamma}}G-\E_{\nu_{N,u,\gamma}}G\bigr\rvert\le2\norm G_\infty\,\mu_{N,\gamma}(S_{N,u,\gamma}^c).
\end{equation}
If $u=nB^r$, the sequence of sides $uB^K$ is the sequence $nB^K$ shifted by $r$. Then the first stage shows that, under $\mu_{nB^K,\gamma}$, the expectations of a fixed bounded continuous function of finitely many flowed Wilson loops form a Cauchy sequence up to an error $4\norm G_\infty F(u)$. Letting $u\to\infty$, we obtain convergence under the unconditioned Wilson measure.
\end{enumerate}
This order of limits, first $K\to\infty$ at fixed $u$ and then $u\to\infty$, appears throughout the paper. Configurations outside the small-field event enter only through the bound on large fields. We prove that bound directly under $\mu_{N,\gamma}$, without conditioning. In particular, we never need to show that a conditioned measure is reflection positive.

The rest of the proof of \cref{s2:thm:main} consists of three tasks. The first is the renormalization group analysis on the small-field event, and its consequences at a fixed number of steps (\cref{s3:rg,s3:fixed}). The second is the bound on large fields (\cref{s3:largefield}). The third is the passage from the renormalization group variables to flowed Wilson loops, and then to a measure on the trajectory space (\cref{s3:flow,s3:stochastic,s3:construction}).

\section{The renormalization group on the small-field event}\label{s3:rg}

\paragraph{Averaging.}
The renormalization group step is built from Bałaban's averaging operation for lattice gauge fields \cite{Balaban1985Averaging}. Fix a configuration on a lattice of side $Mm$. We divide the lattice into cubic blocks of $M^3$ sites, centered at the points of a sublattice of side $m$. For each link $c$ of the coarse lattice, we consider the $M^3$ paths of length $M$ parallel to $c$ that start at the sites of the block at the initial point of $c$. They are joined to the endpoints of $c$ by fixed connecting paths along the coordinate directions. We take the logarithms of their transports relative to the transport along $c$, average them, and multiply the exponential of the average by the transport along $c$; see \eqref{ch01:eq:ordinary-map}. When some relative transport equals $-\Id$, the logarithm is not defined, and we use the reference transport instead. In this way the block averaging map $T_M$ is a Borel map on the whole configuration space. It is gauge covariant, and it commutes with translations of the coarse lattice and with reflections of coordinates that preserve it. (Because the connecting paths change coordinates in a fixed order, we do not claim that it commutes with permutations of coordinates.) Since $M$ is odd, the blocks are centered, and the blocks of successive averaging steps are nested. As a consequence, iterated averages of a configuration supported on one side of a reflection plane depend only on that side (\cref{ch03:chapter}). In \cref{part:finite} we prove the basic properties of the block averaging map, including an exact formula for the integral over a fiber $\{U:T_MU=b\}$ in local coordinates and an estimate that controls the curvature of $T_MU$ by that of $U$ (\cref{ch04:thm:energy-near-contraction}). One renormalization group step consists of $J$ consecutive averaging steps, $V^{(k+1)J}=T_M^JV^{kJ}$. The integer $J$ is fixed at the start of the proof.

\paragraph{Background fields.}
Let $c$ be a coarse configuration close to the identity. The integral over fine configurations $U$ with $T_M^JU=c$ is dominated by a neighborhood of the minimizer of the Wilson action subject to this constraint. Following Bałaban \cite{Balaban1985Variational}, we call this minimizer the background field. We construct it in a gauge-fixed chart near the identity, by solving the constrained Euler--Lagrange equation with a contraction argument (\cref{part:background}). It is a specific branch of local constrained minimizers; we do not claim that it minimizes the action over the whole fiber of the averaging map. We prove that it depends analytically on $c$, with estimates that are uniform in the number of averaging steps between the fine and coarse lattices. The value of the constrained minimum defines the effective classical action $E_d(c)$ at depth $d$, the number of averaging steps. Minimizing in stages gives the nested minimization identity
\[
 E_{d+1}(c)=M\min_zE_d(b(c,z)).
\]
Here $b(c,z)$ parametrizes, in the chart, the configurations on the intermediate lattice whose average is $c$, the minimum is over that chart, and the factor $M$ is the change of units in three dimensions. The effective classical action $E_0$ is the Wilson action. Since $g^2_{d+1}=Mg^2_d$, the ratio $E_d/g_d^2$ is invariant under this change of scale. The main technical point is that the analytic norms of $E_d$ and of the background field do not grow with $d$. This rests on two facts. The first is the linear theory: the effective Maxwell quadratic form has an exponentially decaying kernel, and its restriction to the next constrained fluctuation space has a positive lower bound and an exponentially decaying inverse, uniformly in the depth and the lattice side (\cref{ch05:thm:flat-precision}). Bałaban's propagator estimates \cite{Balaban1984I,Balaban1984II,Balaban1985Background} are used separately, in the constrained linear extension and the scalar gauge inverse of \cref{ch06:thm:matched-flat-extension,ch06:lem:scalar-gauge-inverse}. The second is that the linearized averaging operator, acting on connections after a gauge correction, has row sums one and column sums $M^{-3}$ (\cref{ch07:eq:positive-Q}). As a result the nonlinear errors are summable over scales.

\paragraph{Fluctuations and the effective interaction.}
We fix a renormalization group step from side $M^Jm$ to side $m$. We parametrize the fine configuration by the coarse configuration $c$ and by free logarithmic coordinates $z_i$ at each of the $J$ intermediate scales. We then center these coordinates at the background field and normalize them by the constrained Hessian:
\[
 z_i=z_i^*(c,z_{>i})+g_iK_i(c,z_{>i})^{-1/2}w_i.
\]
Here $z_i^*$ is the conditional constrained minimizer, $K_i$ is the constrained Hessian of the action, and $g_i$ is the coupling at that scale. We call the $w_i$ the normalized fluctuation variables. The change of variables is triangular, from the coarsest intermediate scale to the finest. We compute its Jacobian exactly, including the Haar density and the Jacobian of the constraint (\cref{ch09:chapter}). In the variables $w$ the integrand is $\exp(-V(c,w))$. The interaction $U(c,w)=V(c,w)-\abs w^2/2$ contains the higher-order terms of the action, the Jacobians, and the effective interaction inherited from finer scales. On the small-field event every $w_i$ is bounded by $A_{\mathrm w}\log(1/t)$, where $t$ is the largest coupling of the step. On this product of balls the interaction $U$ is small, and we expand the integral over $w$ in a convergent polymer expansion. We use the cluster expansion in the form of Kotecký and Preiss \cite{KoteckyPreiss1986}, with the tree estimates that go back to Penrose \cite{Penrose1967}, as presented by Fern\'andez and Procacci \cite{FernandezProcacci2007} and Dimock \cite{Dimock2013v2}. After imposing the small-field tests through the $k$-th renormalization group step, we obtain an exact formula, relative to product Haar measure, for the unnormalized stopped density of the averaged configuration:
\begin{equation}\label{s3:eq:effective-density}
 \one_{D_r}(c)\exp\Bigl(\Lambda_k-\frac{E_{kJ,\mathrm{sat}}(c)}{g_{kJ}^2}-\overline R_{kJ,\mathrm{sat}}(c)\Bigr).
\end{equation}
Here $g_{kJ}^2=\gamma/n_{kJ}$ is the coupling of the lattice of side $n_{kJ}$, and $D_r$, the gauge-invariant small-field domain, is the set of configurations that are gauge equivalent to configurations with all link variables close to the identity. The subscript $\mathrm{sat}$ denotes the gauge-invariant extension to $D_r$ of a function defined near the identity. The function $E_{kJ}$ is the effective classical action. The centered effective interaction $\overline R_{kJ}$ is generated by the integration of the fluctuations, and it is normalized so that $\overline R_{kJ}(\Id)=0$. The normalization constant $\Lambda_k$ collects all factors that do not depend on the field, including the vacuum energy (\cref{ch10:chapter,ch11:chapter}). At $k=K$ the normalized stopped density is the law of the coarse field $V^{JK}$ under the conditioned Wilson measure. At intermediate steps it does not include the remaining coarser restrictions of the small-field event.

\paragraph{Analytic norms and contraction.}
We control $\overline R_{kJ}$ by a norm that sums the absolute values of all Taylor coefficients at the identity (\cref{ch01:def:analytic-norms}). The coefficients of order $k$ are weighted by $\rho^{k-2}/k!$ and by $\exp(a_{\mathrm{sp}}\Tree)$, where $\Tree$ is the length of the shortest network that connects the lattice points on which the coefficient depends. The weighted coefficients measure spatial locality. Connected products and compositions are controlled by the coefficient calculus, using output row bounds and volume-normalized input column bounds, under the required image conditions for compositions (\cref{ch10:lem:positive-composition}); derivative estimates hold on strictly smaller analytic domains. Products of interactions whose coefficient networks are not connected, such as products of two extensive interactions, are not controlled in this way. The key estimate is \cref{ch10:thm:common-analytic-class}, which says that one renormalization group step maps a centered effective interaction $R$ to
\[
 R_+=Q+\mathcal P R+\mathcal D R,\qquad \norm{R_+}\le C_G+q\norm R+C_Tt^2\norm R(p^4+\norm Rp^2).
\]
In the notation of that theorem, $Q=Q_{d_0,J}^{\mathrm w}$, $\mathcal P=\mathcal P_\delta$, and $\mathcal D=\mathcal D_J$ for the step under consideration; in this paragraph only, $Q$ does not denote the constant-mode quartic form. The term $Q$ is the newly generated effective interaction, produced by the step itself. The term $\mathcal PR$, the pullback at zero fluctuation, is $R$ evaluated on the background field. It contracts by a factor $q=O(M^{-J})$, because $R$ is gauge invariant, invariant under translations, global color rotations, and inversion, and has no linear term at the identity (\cref{ch08:thm:principal-contraction}). The term $\mathcal DR$ is the fluctuation correction to the pullback. It gains two powers of the coupling, because the linear term of the fluctuation integral vanishes on each connected support. We choose $J$ so that $q<1/4$, and then we choose the coupling small. The ball $\norm R\le E_*$ is then invariant, uniformly in the number of steps. Gauge invariance also enters in another way. At each step we choose a gauge transformation that depends only on the coarse configuration, and that brings the whole reconstructed fine configuration, fluctuations included, into a fixed analytic neighborhood of the identity (\cref{ch10:thm:whole-image-gauge}).

\paragraph{The conditional law of the fluctuations.}
Given the coarse configuration $c$, the normalized fluctuation variables $w$ of one step have a conditional law, the conditional fluctuation kernel, with density proportional to $\one_{\Omega_p}(w)e^{-V_c(w)}$. Here $\Omega_p$ is a product of three-dimensional balls, and $D_w^2V_c\ge\Id/2$ (\cref{ch10:thm:complete-white-law}). We emphasize that this law is not Gaussian, and that the variables are not independent. It is, however, uniformly log-concave on a convex set. Using the Bakry--\'Emery criterion \cite{BakryEmery1985} with reflecting boundary conditions and the Helffer--Sj\"ostrand representation of covariances \cite{HelfferSjostrand1994}, we obtain a logarithmic Sobolev inequality and exponentially decaying covariance estimates, with constants that do not depend on the dimension (\cref{ch02:thm:reflecting-concentration}). These estimates, together with bounds on the conditional means, are the probabilistic input for everything that follows. The conditional means are small, of order $t\norm{\log c}_\infty$; at $c=\Id$ they vanish by the global $\SU$ symmetry (\cref{ch10:thm:actual-means-and-scores}).

\section{Convergence at fixed resolution and reconstruction}\label{s3:fixed}

\paragraph{Convergence of the effective interactions.}
The effective interaction after one step depends on the effective interaction at finer scales only through the contraction above. The influence of the finest scales therefore decays geometrically with the number of steps. In \cref{ch12:chapter} we use this to prove convergence at fixed resolution. We fix the coarse side and a number of steps counted from the coarse side. Then the effective interactions, the conditional laws of the fluctuations, and the joint laws of the coarse configuration and the fluctuations of these steps converge as the total number of steps $K\to\infty$. This is a statement about finitely many variables. A subtle point is that gauge invariance allows different choices of gauge representatives at different $K$. We fix, once and for all, a family of gauge transformations that depend only on the coarse configuration and are chosen by a compactness argument. For this family the reconstructed configurations at each fixed resolution converge (\cref{ch12:thm:replacement-family}).

\paragraph{Reconstruction of the fine field.}
From the coarse configuration and the fluctuations of all steps, we reconstruct the fine configuration step by step, from coarse to fine. We write $A$ for the logarithm of the reconstructed fine configuration in the chosen gauge, divided by the mesh, so that $A$ approximates a connection. In \cref{part:reconstruction} we estimate how $A$ changes when we add the fluctuations of one more step. The change has three parts: a conditionally centered linear part, a nonlinear part, and the difference between the background field and the interpolation of the coarser field. The key deterministic fact concerns the last part. For each input direction and color, the linearized background field has the same signed volume-normalized column integral as the piecewise constant embedding of the coarse field. Their difference therefore has mean zero in each input column, and after heat smoothing at time $s$ it gains a factor $\min(1,a/\sqrt s)$, where $a$ is the coarse mesh (\cref{ch14:chapter}). The key probabilistic fact is concentration for Lipschitz functions of the strongly log-concave law of the fluctuations (\cref{ch13:chapter}). Together, these two facts show that the heat evolution of $A$ is controlled in norms of negative regularity, with summable contributions from the fine scales.

\section{The bound on large fields}\label{s3:largefield}

\Cref{s2:thm:large-field} has two parts. The second term in \eqref{s2:eq:large-field} accounts for configurations whose curvature is small at every scale, but for which some normalized fluctuation variable exceeds its cutoff $A_{\mathrm w}\log(1/t)$. By concentration for the law of the fluctuations, this event has conditional probability at most $\exp(-a\log^2(1/t))$ at each step. The sum over steps converges, because the effective coupling $t$ decreases geometrically toward the ultraviolet. The point that needs care is that we condition only on the preceding scales being good, and never on good behavior at later scales. The probability of the preceding event is multiplied, never divided (\cref{ch19:chapter,ch20:chapter}).

The first term accounts for configurations with a plaquette of large curvature at some scale. Here the argument has three steps.

\paragraph{Step 1: exponential moments of blocked actions.}
Under the fixed-parameter hypotheses of \cref{ch19:thm:all-phase-window}, fix $0<\Gamma<\infty$, take $u\in4\N$ with $u\ge n_{\min}(\Gamma)$, and let $N=uB^K$ and $1\le j\le JK$. Fix a tiling of the lattice of side $n_j$ into cells. Let $Q_j$ be the Wilson action, divided by the coupling, of the averaged configuration $V^{j-1}$ on the plaquettes that determine the curvature of the faces inside the cells. We prove that
\[
 \log\E_{\mu_{N,\gamma}}e^{sQ_j}\le C_vn_{j-1}^3+C_sn_{j-1}^2\log(1/g_{n_{j-1}})+C_b,\qquad0\le s\le s_0,
\]
with constants that do not depend on the number of averaging steps or on $\gamma\in(0,\Gamma]$ (\cref{ch19:thm:all-phase-window}). The surface term comes from the difference between the numbers of independent plaquette variables of the periodic and the open lattice, in a comparison of partition functions, and it has to be kept. This is a stability bound in the spirit of Bałaban, but for an exponential moment under the unconditioned measure. Its proof is the most involved part of the paper. We partition the configurations according to the regions where some averaged plaquette has large curvature, recorded scale by scale; we call this partition the large-field decomposition (\cref{ch18:thm:history-cover}). Outside these regions we run the renormalization group in a localized form. We construct background fields, fluctuation integrals, and effective interactions on arbitrary unions of blocks, with the link variables in the bad regions held fixed as parameters (\cref{ch15:chapter,ch16:chapter,ch17:chapter}). The estimates are uniform in the shape, the number of components, and the age of the bad regions. In \cref{ch18:chapter} we bound the density of the fine configurations on each piece of this partition by the localized effective density with the same normalization. The bound holds pointwise on gauge sections that contain every point of the relevant fibers (\cref{ch18:thm:canonical-density,ch18:thm:slice-induction}). The difference between the localized effective interactions and those on the full torus decays exponentially away from the bad regions. The missing coordinates, the boundaries, and the frozen coordinates contribute to the normalization constant, and we keep track of these contributions separately. Inside the bad regions we use only a crude bound: the effective classical action of the coarse configuration dominates both the Wilson action of the omitted links and a positive cost for each recorded bad plaquette (\cref{ch19:chapter}). This cost grows linearly in $n_j/\gamma$, while the entropy of the bad regions grows only like the square of a logarithm.

\paragraph{Step 2: chessboard estimates.}
The Wilson measure is reflection positive with respect to reflections in lattice planes through sites and through links, as Osterwalder and Seiler showed \cite{OsterwalderSeiler1978}. We use this separately for each of the eight shifted tilings of the lattice of side $n_j$ into cells of $2\times2\times2$ sites. For a cell $C$, let $I_0$ be the event that some internal plaquette of $C$ at level $j$ has large curvature. Let $R_0$ be the event that every plaquette at a finer level whose full dependence set on the fine lattice lies in $C$ has small curvature; plaquettes whose dependence sets cross the boundary of the cell are not tested. Both indicators are bounded, gauge-invariant Borel functions of the fine links in the cell. Because the blocks of the averaging map are centered and nested, the reflections between cells map these dependence sets and tests to their reflected placements in the other cells. Then the chessboard estimate of Fr\"ohlich, Israel, Lieb, and Simon \cite{FILS1978,FILS1980}, in the form given by Biskup \cite{Biskup2009}, applies to this tiling (\cref{ch03:thm:original-chessboard,ch19:lem:marked-chessboard}). Let $R$ and $I_{\mathrm{all}}$ be the intersections of the reflected placements of $R_0$ and of $I_0$ over all cells of the tiling. The chessboard estimate gives
\[
 \mu_{N,\gamma}(R_0\cap I_0)\le\mu_{N,\gamma}(R\cap I_{\mathrm{all}})^{8/n_j^3}.
\]
The regularity events stay inside the root, so no division by their probability occurs. On $R$, the estimate on the curvature of averaged configurations shows that each cell with a plaquette of large curvature contributes at least a fixed multiple of $n_j/\gamma$ to $Q_j$. By Step 1 and Markov's inequality, the probability of $R\cap I_{\mathrm{all}}$ is then at most $\exp(-cn_j^4/\gamma+Cn_j^3+Cn_j^2\log(1/g_{n_j}))$. Taking the root, we find that $R_0\cap I_0$ has probability at most $\exp(-c'n_j/\gamma)$, once $n_j/\gamma$ is large.

\paragraph{Step 3: first failure.}
We decompose the event that some scale has a bad cell according to the first scale, counted from the finest lattice toward the coarse lattice, at which a bad cell occurs. On each piece of this decomposition all finer scales are good, which implies the regularity events $R_0$ of Step 2. We sum the bounds of Step 2 over cells and scales (\cref{ch20:thm:shell-induction}). Since the bound at side $n_j$ decays like $\exp(-cn_j/\gamma)$ and the sides grow geometrically toward the ultraviolet, the sum is dominated by the coarsest side $u$, which gives the first term in \eqref{s2:eq:large-field}. On the complementary event the curvature is small at every scale. By \cref{ch04:thm:small-curvature-admission}, at every scale there exists a gauge representative with all link variables close to the identity. The remaining conditions of the small-field event reduce to the cutoffs on the fluctuations treated above.

The parameters are the integer $J$, the analytic radii, the cutoff amplitude $A_{\mathrm w}$, the widths of the boundary layers in the localized analysis, and the minimum coarse side. We choose them in a fixed order: first the quantities that are uniform in $J$, then $J$, then the constants that depend on $J$, and finally the minimum side (\cref{ch20:thm:common-parameters}). One of the tasks of \cref{ch20:chapter} is to check that this order is free of circularity.

\section{The gradient flow}\label{s3:flow}

To pass from the renormalization group variables to flowed Wilson loops, we need a theory of the lattice gradient flow for rough initial data that is uniform in the mesh. We write a lattice configuration near the identity as $U_e=\exp(aA_e)$. We apply a time-dependent gauge transformation that adds a damping term, in the spirit of the method of Zwanziger \cite{Zwanziger1981}, DeTurck \cite{DeTurck1983}, and Donaldson \cite{Donaldson1985}. Then the Wilson flow becomes a parabolic equation whose linear part is the lattice heat equation,
\[
 \partial_tA=\Delta_aA+\mathcal N_a(A,A)+(\text{higher order}),
\]
where $\mathcal N_a$ is a quadratic nonlinearity involving one lattice derivative. For the initial data of interest, $A$ has regularity just below $-1/2$, and the quadratic term is not defined by naive multiplication. We follow the approach of Da Prato and Debussche \cite{DaPratoDebussche2003}, which Cao and Chatterjee used for the Yang--Mills heat flow with random data \cite{CCHeat2022v4}, and write
\[
 A(t)=x(t)+y(t)+r(t),\qquad x(t)=e^{t\Delta_a}A(0),\qquad y(t)=\int_0^te^{(t-s)\Delta_a}\mathcal N_a(x(s),x(s))\,ds.
\]
We then solve for the remainder $r$, which is more regular. We call the pair $(x,y)$, consisting of the linear heat evolution and the quadratic Picard term, the flow data. For a fixed $\alpha\in(1/2,2/3)$ and horizon $T$, we measure it in the norms
\[
 \sup_{0\le t\le T}\bigl\{(t+a^2)^{\alpha/2}\norm{x(t)}_\infty+(t+a^2)^{(\alpha+1)/2}\norm{D_ax(t)}_\infty\bigr\}
\]
and
\[
 \sup_{0\le t\le T}\bigl\{(t+a^2)^{\alpha-1/2}\norm{y(t)}_\infty+(t+a^2)^{\alpha}\norm{D_ay(t)}_\infty\bigr\},
\]
where $D_a$ includes the lattice differences and shifts that the equation requires (\cref{ch01:def:enhancement}). \Cref{part:flow} proves the deterministic theory. It consists of local existence of $r$ on a short time interval with bounds uniform in the mesh (\cref{ch21:thm:enhanced-local-flow}); smoothing after a positive time; continuation to any fixed time interval, by comparison with finitely many smooth solutions of the continuum flow (\cref{ch21:thm:continuum-atlas,ch22:thm:positive-window-bootstrap}); and Lipschitz dependence of the flowed holonomies on the flow data $(x,y)$ (\cref{ch22:cor:bounded-input-trace-stability}). For smooth initial data the lattice flow converges to the continuum flow. Here we use the deterministic results of Cao and Chatterjee \cite{CCHeat2022v4,CCState2023v4}, which rest on the global theory of R\aa de \cite{Rade1992} and of Charalambous and Gross \cite{CharalambousGross2013}.

\section{Stochastic estimates and convergence of flowed loops}\label{s3:stochastic}

In \cref{part:stochastic} we prove two facts about the flow data $(x,y)$ of the reconstructed fine field under the conditioned Wilson measure. First, it is bounded in the norms of \cref{s3:flow} with high probability, uniformly in the number of steps. Second, it is well approximated by the flow data of the field reconstructed from finitely many steps. At each step the new fluctuations contribute four terms: a conditionally centered linear term, a conditionally centered nonlinear term, an interpolation defect, and a predictable term (\cref{ch23:lem:four-piece-increment}). We estimate all sixteen ordered products of these terms in the quadratic Picard term $y$ (\cref{ch23:chapter,ch24:chapter}). We keep their conditional means, which are nonzero, and we treat both orders of each product, since $\mathcal N_a$ is not symmetric. At mesh $a$ the linear fluctuation contributes $g_3a^{\alpha-1/2}\sqrt{\log(2+L/a)}$, which is summable over the geometric sequence of meshes since $\alpha>1/2$.

Given these estimates, we compare the flowed Wilson loops of the fine configuration with those of the configuration reconstructed from the coarse field and the fluctuations of a fixed number $q$ of steps. They are close in probability, uniformly in the total number of steps. The latter configuration is a smooth interpolation of a finite-dimensional random variable whose law converges as $K\to\infty$, by \cref{s3:fixed}. Further, the flow of a fixed smooth field converges to the continuum flow. Letting first $K\to\infty$ and then $q\to\infty$, we find that the laws of flowed Wilson loops under the conditioned Wilson measure converge (\cref{ch25:thm:selected-loop-limit}). Together with \eqref{s3:eq:transfer} and the bound on large fields, this gives convergence under the Wilson measure (\cref{ch25:thm:original-loop-stitching}). We first prove these results when the coupling of the coarse lattice is below a small ceiling, and then extend them to every fixed $\gamma$ by taking the coarse side large (\cref{ch25:lem:finite-ceiling-selected-chain,ch26:thm:finite-ceiling-matching}). This proves part (1) of \cref{s2:thm:main} along the subsequences $N_K=n_*B^K$.

\section{Construction of the limiting measure}\label{s3:construction}

Convergence of finite families of flowed Wilson loops does not by itself produce a measure on $\OrbitSpace$. We obtain one as follows (\cref{ch26:chapter}). For each fixed number $q$ of steps, the field reconstructed from $q$ steps is smooth after interpolation, and its continuum flow defines a random element of $\OrbitSpace$. These random trajectories have energies that are tight uniformly in $q$ at each of the times $s_j=L^22^{-j}$. Indeed, the lattice energy along the flow is controlled by the stochastic estimates, and it does not increase in time. The sets of trajectories whose energies at the times $s_j$ are bounded by given numbers $R_j$, for all $j$, are compact, by results of Cao and Chatterjee that rest on Uhlenbeck's compactness theorem. Hence these laws converge along subsequences, and the joint laws of flowed Wilson loops identify the limit. The joint convergence of energies and Wilson loops follows from a comparison with the smooth fields reconstructed from finitely many steps, together with the continuity of the energy after a positive time in the Sobolev setting of Cao and Chatterjee (\cref{ch26:thm:selected-energy-loop-limit}).

We also want to realize the limit as a limit of the lattice measures themselves. For this purpose we construct measurable gauge-invariant maps $\pi_K$ from lattice configurations to $\OrbitSpace$. We enumerate the trajectories of the continuum flow started from trigonometric polynomials with rational coefficients. Given a lattice configuration, $\pi_K$ selects the first trajectory in this list with the following two properties: its first $m$ determining coordinates are within $2^{-m}$ of the flowed lattice Wilson loops, and its energies at $m$ fixed positive times are at most the lattice energies plus $2^{-m}$. If there is no such trajectory, $\pi_K$ selects the zero trajectory. With $m=m(K)$ increasing slowly, these maps have the properties in part (3) of \cref{s2:thm:main}, and the energy constraint gives the tightness in part (4) (\cref{ch26:thm:energy-constrained-matching,ch26:thm:orbit-limit}).

Finally, the passage from the subsequences $n_*B^K$ to the full sequence $nM^k$ is arithmetic. Since the construction works for every sufficiently large $J$, we may take two consecutive values $J$ and $J+1$. Every residue class of $k$ modulo $J$ contains infinitely many $k$ for which $nM^k$ also belongs to a sequence of the form $n'M^{(J+1)K}$, and hence the limits along the different residue classes coincide (\cref{ch26:thm:full-M-ladder}). We prove the remaining parts of \cref{s2:thm:further} in \cref{ch26:chapter-properties}.

\section{Short distances}\label{s3:short}

For the short-distance theorem we fix $\gamma$ and a coarse side. We study the heat-smoothed circulation of the reconstructed field around a small loop $x+\varepsilon\ell$, normalized by $\sqrt{g_3^2\varepsilon}$. Adding the fluctuations step by step, we write this circulation as the sum of three terms: a contribution from the coarse field, a martingale whose increments are linear in the new normalized fluctuations, and a small nonlinear error (\cref{ch28:chapter}). A useful observation is that the linearized successive background fields differ from a single background field by a periodic gradient, which the closed circulation does not see. We compare each martingale increment to a Gaussian, by an integration by parts adapted to the cutoff on the fluctuations, in the spirit of Stein's method \cite{Stein1972NormalApproximation}. The sum of the conditional covariances converges to the covariance \eqref{s2:eq:circulation-covariance} of the free field, as $\varepsilon\to0$ after the lattice spacing (\cref{ch28:thm:maxwell-calibration,ch29:chapter}). The contribution of the coarse field vanishes in the limit because the loop shrinks. We then compare the flowed Wilson loop with the circulation, using the stability of the flow and the expansion $1-\tr e^V/2=\abs V^2/8+O(\abs V^4)$ (\cref{ch29:lem:actual-flow-trace-reduction}). We carry out this analysis under the conditioned Wilson measure and in the fixed gauge of the reconstruction. To obtain the gauge-invariant statement under $\lambda$, we take the limits in the following order: the lattice spacing to zero; then $\varepsilon\to0$ at a fixed coarse side and with a fixed cutoff on the flow data; then the removal of that cutoff; and finally the coarse side to infinity. The asymptotic independence from macroscopic observables has a simple reason: the Gaussian approximation holds uniformly conditionally on any fixed number of coarse steps, while a macroscopic flowed observable is close to a function of finitely many coarse steps.

\section{Weak coupling}\label{s3:weak}

For the weak-coupling results we fix, before $\gamma$ varies, the parameters of the construction and a coarse side $u$ in the class $[n]_M$, larger than certain thresholds; by \cref{s2:thm:main}(5) this does not change $\lambda_\gamma$. As $\gamma\to0$ at this fixed coarse side, every scale of the hierarchy becomes weakly coupled, and the analysis takes place on the coarse lattice. \Cref{ch27:chapter} develops the geometry near flat connections and Laplace's method on the coarse lattice. The density of the coarse configuration at the end of the renormalization group is proportional to $\one_{D_r}\exp(-uE_d/\gamma-R+R(\Id))$, and its normalizing constant is comparable to $\gamma^{3u^3-1}$, uniformly in the number of finer steps (\cref{ch27:thm:sharp-normalizer}). Near a regular flat connection, the integral over the directions transverse to the flat connections is Gaussian. In \cref{ch30:chapter} we compare the reconstructed field with the field in which all fluctuations beyond a fixed number of steps are set to zero. Under the conditioned Wilson measure, their flowed holonomies differ by $O(\sqrt\gamma)$ in $L^p$, uniformly in the number of steps and in small $\gamma$ (\cref{ch30:thm:paired-enhancement-return,ch30:thm:scaled-matrix-tail}). This gives the concentration on flat connections in \cref{s2:thm:flat}.

For \cref{s2:thm:regular}, recall that the normal covariance $\NormalCov$ is the covariance of the linear response of the flowed logarithms at a regular flat connection. We compute it explicitly in \cref{ch27:thm:flowed-log-covariance}. The theorem then follows from Gaussian integration at finitely many scales and the transfer estimates of \cref{ch30:chapter} (\cref{ch31:chapter}). We take the limits in the following order: the continuum limit, then $\gamma\to0$ at a fixed number of scales, and then the number of scales to infinity.

For \cref{s2:thm:central}, the starting point is an exact Taylor expansion of the effective classical action of a constant configuration, $E_d(e^{X/u})=Q(X)/u+O(\abs X^5)$, where $Q$ is the constant-mode quartic form and the remainder is uniform in the depth $d$ (\cref{ch32:thm:absolute-quartic}). We integrate out the nonconstant modes with Gaussian estimates. This leaves a nine-dimensional integral against $e^{-Q(Y)}$ on the scale $Y=X/\gamma^{1/4}$, multiplied by a positive factor that is bounded above and below (\cref{ch32:thm:common-scalar-core}). This factor is the same for all test functions, so it cancels in the ratios \eqref{s2:eq:core-ratio}. The expectation of a compactly supported test function at this scale is of order $\gamma^{1/4}$, which comes from the volume $\gamma^{9/4}$ of the constant modes, compared with the order $\gamma^2$ of the normalization. We then transfer the result from the reconstructed holonomies to the flowed holonomies of the field. For this we use an estimate that is conditional on the coarse configuration (\cref{ch30:thm:terminal-frozen-matrix-tail}). An estimate averaged over the coarse configuration would lose the factor $\gamma^{1/4}$ when divided by this small expectation.

\section{Order of choices and dependence of constants}\label{s3:constants}

Many statements in the paper are uniform in some parameters and not in others, and the order in which parameters are chosen matters. We describe the convention used throughout in \cref{ch01:conv:constants}. In brief, we choose the integer $J$ first, using only estimates that are uniform in $J$. Then we compute all local constants at that $J$. They are uniform in the lattice side, the number of averaging steps, the coupling below the relevant ceiling, and, in the localized analysis, the shape of the regions. Finally we choose the minimum coarse side. Statements at fixed $\gamma$ may have constants that depend on $\gamma$; these include the convergence of flowed loops, the construction of $\lambda$, and the short-distance theorem. The bound on large fields and the weak-coupling estimates are uniform in $\gamma$ below fixed ceilings.
